\frfilename{mt1a3.tex}
\versiondate{18.12.03}
\copyrightdate{1994}

\def\chaptername{Lampiran}
\def\sectionname{Limit superior dan limit inferior}

\newsection{1A3}

Tidak semua mata kuliah pengantar analisis real memiliki cukup waktu untuk
membahas konsep $\limsup$ dan $\liminf$.

\leader{1A3A}{Definisi (a)} Untuk suatu barisan bilangan real $\sequencen{a_n}$,
tuliskan

\Centerline{$\limsup_{n\to\infty}a_n=\lim_{n\to\infty}\sup_{m\ge n}a_m
=\inf_{n\in\Bbb N}\sup_{m\ge n}a_m$,}

\Centerline{$\liminf_{n\to\infty}a_n=\lim_{n\to\infty}\inf_{m\ge n}a_m
=\sup_{n\in\Bbb N}\inf_{m\ge n}a_m$;}

\noindent jika kita mengizinkan nilai $\pm \infty$, baik ketika membentuk
supremum dan infimum maupun ketika mengambil limit\cmmnt{ (lihat 112Ba)},
keduanya akan selalu terdefinisi\cmmnt{, karena barisan-barisan

\Centerline{$\sequencen{\sup_{m\ge n}a_m}$,
\quad$\sequencen{\inf_{m\ge n}a_m}$}

\noindent bersifat monoton}.


\cmmnt{\header{1A3Ab}{\bf (b)} Secara eksplisit:


\Centerline{$\limsup_{n\to\infty}a_n=\infty\,\iff\,\{a_n:n\in\Bbb N\}$
tidak terbatas di atas,}

\Centerline{$\limsup_{n\to\infty}a_n=-\infty\,\iff\,\lim_{n\to\infty}a_n
=-\infty$,}

\noindent yakni, jika dan hanya jika untuk setiap $a\in\Bbb R$ terdapat
$n_0\in\Bbb N$ sedemikian sehingga $a_n\le a$ untuk setiap $n\ge n_0$;

\Centerline{$\liminf_{n\to\infty}a_n=-\infty\, \iff\,\{a_n:n\in\Bbb N\}$
tidak terbatas di bawah,}

\Centerline{$\liminf_{n\to\infty}a_n=\infty\, \iff
\,\lim_{n\to\infty}a_n
=\infty$,}

\noindent yakni, jika dan hanya jika untuk setiap $a\in\Bbb R$ terdapat
$n_0\in\Bbb N$ sedemikian sehingga $a_n\ge a$ untuk setiap $n\ge n_0$.
}%end of comment

\header{1A3Ac}{\bf (c)} Untuk \cmmnt{$a\in\Bbb R$ berhingga, kita mempunyai

\inset{$\limsup_{n\to\infty}a_n=a$ jika dan hanya jika (i) untuk
setiap $\epsilon>0$ terdapat $n_0\in\Bbb N$ sedemikian sehingga $a_n\le
a+\epsilon$ untuk setiap $n\ge n_0$ (ii) untuk setiap $\epsilon>0$ dan setiap
$n_0\in\Bbb N$ terdapat $n\ge n_0$ sedemikian sehingga $a_n\ge a-\epsilon$,}

\noindent sedangkan

\inset{$\liminf_{n\to\infty}a_n=a$ jika dan hanya jika (i) untuk
setiap $\epsilon>0$ terdapat $n_0\in\Bbb N$ sedemikian sehingga $a_n\ge
a-\epsilon$ untuk setiap $n\ge n_0$ (ii) untuk setiap $\epsilon>0$ dan setiap
$n_0\in\Bbb N$ terdapat $n\ge n_0$ sedemikian sehingga $a_n\le a+\epsilon$.}

\noindent Secara umum, untuk }%end of comment
$u\in[-\infty,\infty]$, kita dapat mengatakan bahwa

\inset{$\limsup_{n\to\infty}a_n=u$ jika dan hanya jika (i) untuk setiap
$v>u$ (jika memang ada) terdapat $n_0\in\Bbb N$ sedemikian sehingga $a_n\le v$
untuk setiap $n\ge n_0$ (ii) untuk setiap $v<u$ dan setiap $n_0\in\Bbb N$
terdapat $n\ge n_0$ sedemikian sehingga $a_n\ge v$,}

\inset{$\liminf_{n\to\infty}a_n=u$ jika dan hanya jika (i) untuk setiap
$v<u$ terdapat $n_0\in\Bbb N$ sedemikian sehingga $a_n\ge v$ untuk setiap
$n\ge n_0$ (ii) untuk setiap $v>u$ dan setiap $n_0\in\Bbb N$ terdapat
$n\ge n_0$ sedemikian sehingga $a_n\le v$.}


\leader{1A3B}{}\cmmnt{ Kita mempunyai hasil-hasil dasar berikut.

\medskip

\noindent}{\bf Proposisi} Untuk sembarang barisan $\sequencen{a_n}$,
$\sequencen{b_n}$ dalam $\Bbb R$,

(a) $\liminf_{n\to\infty}a_n\le\limsup_{n\to\infty}a_n$,

(b) $\lim_{n\to\infty}a_n=u\in[-\infty,\infty]$ jika dan hanya jika
$\limsup_{n\to\infty}a_n=\liminf_{n\to\infty}a_n=u$,

(c) $\liminf_{n\to\infty}a_n=-\limsup_{n\to\infty}(-a_n)$,

(d) $\limsup_{n\to\infty}(a_n+b_n)
\le\limsup_{n\to\infty}a_n+\limsup_{n\to\infty}b_n$,

(e) $\liminf_{n\to\infty}(a_n+b_n)
\ge\liminf_{n\to\infty}a_n+\liminf_{n\to\infty}b_n$,

(f) $\limsup_{n\to\infty}ca_n=c\limsup_{n\to\infty}a_n$ jika $c\ge 0$,

(g) $\liminf_{n\to\infty}ca_n=c\liminf_{n\to\infty}a_n$ jika $c\ge 0$,

\noindent dengan ketentuan bahwa dalam (d) dan (e) kita harus dapat
menafsirkan ruas kanan pertidaksamaan menurut aturan dalam 135A, sedangkan
dalam (f) dan (g) kita mengambil $0\cdot\infty=0\cdot(-\infty)=0$.

\proof{{\bf (a)} $\sup_{m\ge n}a_m\ge\inf_{m\ge n}a_m$ untuk setiap
$n$, sehingga

\Centerline{$\limsup_{n\to\infty}a_n=\lim_{n\to\infty}\sup_{m\ge
n}a_m\ge\lim_{n\to\infty}\inf_{m\ge n}a_m=\liminf_{n\to\infty}a_n$.}

\medskip

{\bf (b)} Dengan menggunakan uraian terakhir mengenai
$\limsup_{n\to\infty}$ dan $\liminf_{n\to\infty}$ dalam 1A3Ac, serta
uraian yang bersesuaian mengenai $\lim_{n\to\infty}$, kita mempunyai

$$\eqalign{\lim_{n\to\infty}&a_n=u\cr
&\iff\text{ untuk setiap }v>u\text{ terdapat }n_1\in\Bbb N
    \text{ sedemikian sehingga }a_n\le v\text{ untuk setiap }n\ge n_1\cr
&\qquad\text{ dan untuk setiap }v<u\text{ terdapat }n_2\in\Bbb N
    \text{ sedemikian sehingga }a_n\ge v\text{ untuk setiap }n\ge n_2\cr
&\iff\text{ untuk setiap }v>u\text{ terdapat }n_1\in\Bbb N
    \text{ sedemikian sehingga }a_n\le v\text{ untuk setiap }n\ge n_1\cr
  &\qquad\text{ dan untuk setiap }v<u,\,\text{serta setiap }n_0\in\Bbb N
    \text{ terdapat }n\ge n_0\text{ sedemikian sehingga }a_n\ge v\cr
&\qquad\text{ dan untuk setiap }v<u\text{ terdapat }n_2\in\Bbb N
    \text{ sedemikian sehingga }a_n\ge v\text{ untuk setiap }n\ge n_2\cr
  &\qquad\text{ dan untuk setiap }v>u,\,\text{serta setiap }n_0\in\Bbb N
    \text{ terdapat }n\ge n_0\text{ sedemikian sehingga }a_n\le v\cr
&\iff\limsup_{n\to\infty}a_n=\liminf_{n\to\infty}a_n=u.\cr}$$

\medskip

{\bf (c)} Ini sekadar soal membalik tanda dan menukar supremum dengan
infimum dalam rumus-rumus tersebut:

$$\eqalign{\liminf_{n\to\infty}a_n
&=\sup_{n\in\Bbb N}\inf_{\lower 0.4ex\hbox{$\scriptstyle m\ge n$}}a_m
=\sup_{n\in\Bbb N}(-\sup_{m\ge n}(-a_m))\cr
&=-\inf_{\lower 0.4ex\hbox{$\scriptstyle n\in\Bbb N$}}
     \sup_{m\ge n}(-a_m)
=-\limsup_{n\to\infty}(-a_n).\cr}$$

\medskip

{\bf (d)} Jika $v>\limsup_{n\to\infty}a_n+\limsup_{n\to\infty}b_n$,
terdapat $v_1$, $v_2$ sedemikian sehingga
$v_1>\limsup_{n\to\infty}a_n$, $v_2>\limsup_{n\to\infty}b_n$ dan
$v_1+v_2=v$. Sekarang terdapat $n_1$, $n_2\in\Bbb N$ sedemikian sehingga
$\sup_{m\ge n_1}a_m\le v_1$ dan $\sup_{m\ge n_2}b_m\le v_2$; sehingga

$$\eqalign{\sup_{m\ge\max(n_1,n_2)}(a_m+b_m)
&\le\sup_{m\ge\max(n_1,n_2)}a_m+\sup_{m\ge\max(n_1,n_2)}b_m\cr
&\le\sup_{m\ge n_1}a_m+\sup_{m\ge n_2}b_m
\le v_1+v_2
= v.\cr}$$

\noindent Karena $v$ sembarang,

\Centerline{$\limsup_{n\to\infty}(a_n+b_n)=\inf_{n\in\Bbb N}\sup_{m\ge
n}(a_m+b_m)\le\limsup_{n\to\infty}a_n+\limsup_{n\to\infty}b_n$.}

\medskip

{\bf (e)} Dengan menggabungkan (c) dan (d),

$$\eqalign{\liminf_{n\to\infty}(a_n+b_n)
&=-\limsup_{n\to\infty}((-a_n)+(-b_n))\cr
&\ge-\limsup_{n\to\infty}(-a_n)-\limsup_{n\to\infty}(-b_n)
=\liminf_{n\to\infty}a_n+\liminf_{n\to\infty}b_n.\cr}$$

\medskip

{\bf (f)} Karena $c\ge 0$,

$$\eqalign{\limsup_{n\to\infty}ca_n
&=\inf_{n\in\Bbb N}\sup_{m\ge n}ca_m
=\inf_{n\in\Bbb N}c\sup_{m\ge n}a_m\cr
&=c\inf_{n\in\Bbb N}\sup_{m\ge n}a_m
=c\limsup_{n\to\infty}a_n.\cr}$$

\medskip

{\bf (g)} Terakhir,

\Centerline{$\liminf_{n\to\infty}ca_n
=-\limsup_{n\to\infty}c(-a_n)
=-c\limsup_{n\to\infty}(-a_n)
=c\liminf_{n\to\infty}a_n$.}
}%end of proof of 1A3B

\cmmnt{
\leader{1A3C}{Catatan} Tentu saja hasil-hasil yang lazim, yakni bahwa
$\lim_{n\to\infty}(a_n+b_n)=\lim_{n\to\infty}a_n+\lim_{n\to\infty}b_n$,
$\lim_{n\to\infty}ca_n=c\lim_{n\to\infty}a_n$, merupakan akibat langsung
dari 1A3B.
}%end of comment

\leader{*1A3D}{Bentuk lain dari gagasan yang sama}
Konsep $\limsup$ dan $\liminf$ dapat diterapkan dalam konteks apa pun yang
memungkinkan kita meninjau limit fungsi bernilai real. Sebagai contoh, jika $f$
adalah fungsi bernilai real yang didefinisikan (sekurang-kurangnya) pada
suatu interval berlubang berbentuk $\{x:0<|c-x|\le\epsilon\}$, dengan
$c\in\Bbb R$ dan $\epsilon>0$, maka

\Centerline{$\limsup_{t\to c}f(t)
=\lim_{\delta\downarrow 0}\sup_{0<|t-c|\le\delta}f(t)
=\inf_{0<\delta\le\epsilon}\sup_{0<|t-c|\le\delta}f(t)$,}

\Centerline{$\liminf_{t\to c}f(t)
=\lim_{\delta\downarrow 0}\inf_{0<|t-c|\le\delta}f(t)
=\sup_{0<\delta\le\epsilon}\inf_{0<|t-c|\le\delta}f(t)$,}

\noindent dengan mengizinkan $\infty$ dan $-\infty$ sesuai kebutuhan.
Atau, jika $f$ didefinisikan pada interval setengah-terbuka
$\ocint{c,c+\epsilon}$, kita dapat mengatakan

\Centerline{$\limsup_{t\downarrow c}f(t)
=\lim_{\delta\downarrow 0}\sup_{c<t\le c+\delta}f(t)
=\inf_{0<\delta\le\epsilon}\sup_{c<t\le c+\delta}f(t)$,}

\Centerline{$\liminf_{t\downarrow c}f(t)
=\lim_{\delta\downarrow 0}\inf_{c<t\le c+\delta}f(t)
=\sup_{0<\delta\le\epsilon}\inf_{c<t\le c+\delta}f(t)$.}

\noindent Demikian pula, jika $f$ didefinisikan pada $\coint{M,\infty}$
untuk suatu $M\in\Bbb R$, kita mempunyai

\Centerline{$\limsup_{t\to\infty}f(t)
=\lim_{a\to\infty}\sup_{t\ge a}f(t)
=\inf_{a\ge M}\sup_{t\ge a}f(t)$,}

\Centerline{$\liminf_{t\to\infty}f(t)
=\lim_{a\to\infty}\inf_{t\ge a}f(t)
=\sup_{a\ge M}\inf_{t\ge a}f(t)$.}

\cmmnt{\noindent Perluasan lebih lanjut dari gagasan ini dibahas secara
singkat dalam 2A3S di Jilid 2.}

\frnewpage
